\documentclass[11pt]{article}
\usepackage[paperwidth=220mm,paperheight=320mm,left=16mm,right=16mm,top=18mm,bottom=18mm,headheight=16pt,headsep=10pt,footskip=14pt]{geometry}
\usepackage{fix-cm}
\usepackage{fontspec,amsmath,amssymb,graphicx,array,multirow,tikz}
\usetikzlibrary{decorations.pathreplacing}
\usepackage[unicode,hidelinks]{hyperref}
\usepackage{polyglossia}
\setmainlanguage{latin}
\setotherlanguages{arabic,greek,hebrew,syriac}
\setmainfont{Linux Libertine O}[Ligatures=TeX]
\setsansfont{Linux Biolinum O}[Ligatures=TeX]
\newfontfamily\greekfont{FreeSerif}[Script=Greek]
\newfontfamily\arabicfont{Amiri}[Script=Arabic]
\newfontfamily\hebrewfont{Frank Ruhl Hofshi}[Script=Hebrew]
\newfontfamily\syriacfont{Segoe UI Historic}[Script=Syriac]
\newfontfamily\CapFont{Noto Serif}
\newfontfamily\MoonFont{FreeSerif}
\newcommand{\MoonSym}{\text{\MoonFont ☾}}\newcommand{\SunSym}{\text{\MoonFont ⊙}}
\newfontfamily\EthiopicFont{Ebrima}
\newcommand{\textethiopic}[1]{{\EthiopicFont #1}}
\newfontfamily\NallinoSigns{NallinoSigns.otf}[Path=./fonts/]
\newcommand{\AbjadZero}{{\NallinoSigns\symbol{"E000}}}
\hypersetup{pdftitle={Al-Battani, Opus astronomicum. Nallino, Index geographicus et Index historicus (Pars II, pp. 359-413). S09},pdfauthor={Al-Battani; Carlo Alfonso Nallino (indices)}}
\makeatletter
\def\ps@sourceedition{%
\def\@oddhead{\vbox{\hbox to\textwidth{\fontsize{8}{10}\selectfont C. A. NALLINO — \IndexHead\hfil\leftmark}\vskip3pt\hrule height.25pt}}%
\let\@evenhead\@oddhead
\def\@oddfoot{\hbox to\textwidth{\fontsize{7}{9}\selectfont\rightmark\hfil\thepage}}%
\let\@evenfoot\@oddfoot}
\makeatother
\newcommand{\IndexHead}{INDEX GEOGRAPHICUS}
\pagestyle{sourceedition}
\setlength{\parindent}{0pt}\setlength{\parskip}{0pt}
\newcommand{\BeginSource}[2]{\clearpage\markboth{PAG. #2}{#1}\hypertarget{#1}{}\pdfbookmark[0]{#1 — p. #2}{bk-#1}\fontsize{11.1}{13.8}\selectfont}
\tracinglostchars=2
\newcommand{\Name}[1]{{\addfontfeatures{LetterSpace=4.0}#1}}
\newcommand{\Indent}{\hspace*{1.6em}}
% an Arabic run inside a Latin line takes no height or depth: the lines are set at their printed baselines
% (@baselines), which the vowel signs of the Arabic open unevenly
\newcommand{\ArabicRun}[1]{\smash{\textarabic{#1}}}
% letters of a geometrical figure, overlined as printed (the line stands 9.3-9.7 pt above the baseline)
\newcommand{\ArOver}[1]{\vbox{\hrule height .4pt\kern 1.2pt\hbox{\rule{0pt}{8.1pt}#1}}}
\newcommand{\qfrac}[2]{{}^{#1}\!/_{{#2}}}
\newcommand{\sa}{\textsuperscript{\textit{a}}}\newcommand{\sm}{\textsuperscript{\textit{m}}}\newcommand{\sd}{\textsuperscript{\textit{d}}}
\newcommand{\RawBlock}[1]{\par\vspace{3pt}\noindent#1\par\vspace{3pt}}
\newcommand{\CenterLine}[1]{\makebox[\linewidth][c]{#1}}
\newlength{\FullSourceWidth}\newlength{\GutterOffset}\newif\ifNumbersLeft
\newcommand{\DSource}[2]{\BeginSource{AB01-PDF#1}{#2}%
 \setlength{\FullSourceWidth}{\textwidth}\setlength{\GutterOffset}{0pt}%
 \ifodd#2\relax\NumbersLefttrue\else\NumbersLeftfalse\fi}
\newsavebox{\DLBox}
\newcommand{\DLine}[3]{%
 \par\noindent\hypertarget{#1}{}%
 \ifx\relax#2\relax\else
   \ifNumbersLeft
    \rlap{\kern-\GutterOffset\llap{\fontsize{7}{8}\selectfont #2\hspace{3mm}}}%
   \else
    \rlap{\kern\dimexpr\FullSourceWidth-\GutterOffset\relax\hspace{3mm}{\fontsize{7}{8}\selectfont #2}}%
   \fi
 \fi
 \sbox{\DLBox}{\strut #3}%
 \ifdim\wd\DLBox>\linewidth\typeout{DIPLOMATICWIDTH|#1|\the\wd\DLBox|\the\linewidth}\fi
 \usebox{\DLBox}\strut\par}
\begin{document}
\thispagestyle{empty}
\begin{center}\Large AL-BATTĀNĪ\par\vspace{4mm}\LARGE OPUS ASTRONOMICUM\par
\vspace{5mm}\large Caroli Alphonsi Nallino index geographicus et index historicus (pars II)\par\vspace{8mm}
\Large S09 --- editio diplomatica\par\vspace{4mm}\large Paginae impressae 359--359\end{center}
\vspace{12mm}\noindent Singulae lineae paginarum impressarum (et columnarum) singulis lineis huius editionis respondent, cum
divisionibus vocabulorum, signis et generibus litterarum (rectis, inclinatis, distantibus). Textus ex imaginibus
paginarum transcriptus est; litterae Graecae et Arabicae, transcriptiones nominum, numeri et signa in imaginibus amplificatis lecta sunt.
Haec transcriptio a nullo homine recognita est.\par
\renewcommand{\IndexHead}{INDEX GEOGRAPHICUS}
\DSource{0808}{359}
\par\noindent\rule{\linewidth}{1.2pt}\par\nointerlineskip\vspace{1pt}\noindent\rule{\linewidth}{.4pt}\par
\par\vspace{32.56mm}
\DLine{AB01-PDF0808-body-L001}{}{\CenterLine{\fontsize{24}{28}\selectfont INDEX GEOGRAPHICUS}}
\par\vspace{6.32mm}
\par\noindent\makebox[\linewidth][c]{\rule[.5ex]{22.44mm}{.4pt}}\par
\par\vspace{12.37mm}
\par\fontsize{9.2}{10.84}\selectfont
\hypertarget{AB01-PDF0808-P01}{}
\DLine{AB01-PDF0808-body-L002}{}{\Indent Numeri inclinatis characteribus impressi indicant Albatenii locos, reliqui adnotationes meas.}
\par\vspace{0.06pt}
\hypertarget{AB01-PDF0808-P02}{}
\DLine{AB01-PDF0808-body-L003}{}{\Indent Compendia haec sunt: lin. = linea. n. = nota, in imis paginis. nr. = numerus. add. = cum}
\par\vspace{-0.09pt}
\DLine{AB01-PDF0808-body-L004}{}{additionibus. add. t. II = cum additionibus initio t. II collocatis. emend. = cum emendationibus.}
\par\vspace{0.09pt}
\hypertarget{AB01-PDF0808-P03}{}
\DLine{AB01-PDF0808-body-L005}{5}{\Indent Permulto post tabulas impressas prodiit optimus liber G. \Name{Le Strange}, \textit{The lands of the Ea-}}
\par\vspace{0.12pt}
\DLine{AB01-PDF0808-body-L006}{}{\textit{stern Caliphate, Mesopotamia, Persia, and Central Asia from the Moslem conquest to the time of Ti-}}
\par\vspace{-0.18pt}
\DLine{AB01-PDF0808-body-L007}{}{\textit{mur}, Cambridge 1905, qui pro compluribus nominibus conferri nunc potest.}
\par\vspace{18.43pt}
\par\nointerlineskip
\noindent\begin{minipage}[t]{.485\textwidth}\vspace{0pt}\vspace*{0.15pt}\fontsize{11.1}{14.02}\selectfont\setlength{\GutterOffset}{0pt}\setlength{\FullSourceWidth}{\linewidth}
\hypertarget{AB01-PDF0808-col_left-E01}{}
\DLine{AB01-PDF0808-col_left-L001}{}{Ābaskūn (Abuskūn, Abiskūn) 172 et n. (et add.}
\par\vspace{0.10pt}
\DLine{AB01-PDF0808-col_left-L002}{}{\hspace*{10.04pt}t. II).}
\par\vspace{0.00pt}
\hypertarget{AB01-PDF0808-col_left-E02}{}
\DLine{AB01-PDF0808-col_left-L003}{}{‘Abbādān II, \textit{53} (et 218).}
\par\vspace{0.10pt}
\hypertarget{AB01-PDF0808-col_left-E03}{}
\DLine{AB01-PDF0808-col_left-L004}{}{Achaia II, \textit{34}.}
\par\vspace{-0.09pt}
\hypertarget{AB01-PDF0808-col_left-E04}{}
\DLine{AB01-PDF0808-col_left-L005}{}{Adana v. Adhanah.}
\par\vspace{0.10pt}
\hypertarget{AB01-PDF0808-col_left-E05}{}
\DLine{AB01-PDF0808-col_left-L006}{}{‘Aden \textit{15}, 170 (et add. t. II); II, \textit{44} (et 217)}
\par\vspace{0.00pt}
\DLine{AB01-PDF0808-col_left-L007}{}{\hspace*{10.04pt}Sinus ‘Aden v. Barbaricus sinus.}
\par\vspace{0.01pt}
\hypertarget{AB01-PDF0808-col_left-E06}{}
\DLine{AB01-PDF0808-col_left-L008}{}{Adhanah II, \textit{39} (et 216: Edane), 43 n. 179.}
\par\vspace{0.18pt}
\hypertarget{AB01-PDF0808-col_left-E07}{}
\DLine{AB01-PDF0808-col_left-L009}{}{Ādharbayǵān II, \textit{35}, 193, 53 n. 266.}
\par\vspace{-0.09pt}
\hypertarget{AB01-PDF0808-col_left-E08}{}
\DLine{AB01-PDF0808-col_left-L010}{}{Adrias v. Sinus Adrias.}
\par\vspace{0.10pt}
\hypertarget{AB01-PDF0808-col_left-E09}{}
\DLine{AB01-PDF0808-col_left-L011}{}{Aegyptus (\textit{Miṣr}), \textit{18}, \textit{19} quater, 172 (et add.}
\par\vspace{0.01pt}
\DLine{AB01-PDF0808-col_left-L012}{}{\hspace*{10.04pt}t. II); II, \textit{34} nr. 42 ac 45; \textit{Eghifṭos} inferior}
\par\vspace{-0.09pt}
\DLine{AB01-PDF0808-col_left-L013}{}{\hspace*{10.04pt}II, \textit{34}, 210. V. Mare Aegypti.}
\par\vspace{-0.09pt}
\hypertarget{AB01-PDF0808-col_left-E10}{}
\DLine{AB01-PDF0808-col_left-L014}{}{Aelana 17 n. 6. V. Aylah.}
\par\vspace{0.01pt}
\hypertarget{AB01-PDF0808-col_left-E11}{}
\DLine{AB01-PDF0808-col_left-L015}{}{Aethiopes, Aethiopia (\textit{al-ḥabash}) \textit{17} ter, \textit{18} bis,}
\par\vspace{0.19pt}
\DLine{AB01-PDF0808-col_left-L016}{}{\hspace*{10.04pt}\textit{19} bis, 170 ter, 171 (et add. t. II); II, \textit{44}.}
\par\vspace{-0.18pt}
\DLine{AB01-PDF0808-col_left-L017}{}{\hspace*{10.04pt}V. Kūsh.}
\par\vspace{0.18pt}
\hypertarget{AB01-PDF0808-col_left-E12}{}
\DLine{AB01-PDF0808-col_left-L018}{}{al-Aflāǵ regio II, 44 n. 192.}
\par\vspace{0.19pt}
\hypertarget{AB01-PDF0808-col_left-E13}{}
\DLine{AB01-PDF0808-col_left-L019}{}{Africa, quo modo eius littus in Albatenii fonte}
\par\vspace{0.00pt}
\DLine{AB01-PDF0808-col_left-L020}{}{\hspace*{10.04pt}descriptum 167. \textit{Ifrīqiyyah}, sensu strictiore}
\par\vspace{-0.17pt}
\DLine{AB01-PDF0808-col_left-L021}{}{\hspace*{10.04pt}Romano, 172; II, \textit{34} (et 215), 212; eius}
\par\vspace{0.10pt}
\DLine{AB01-PDF0808-col_left-L022}{}{\hspace*{10.04pt}mare I, 171. \textit{Ifrīqiyyah}, sensu lato, II, \textit{34}}
\par\vspace{-0.45pt}
\DLine{AB01-PDF0808-col_left-L023}{}{\hspace*{10.04pt}nr. 44. V. Libya.}
\par\vspace{0.28pt}
\hypertarget{AB01-PDF0808-col_left-E14}{}
\DLine{AB01-PDF0808-col_left-L024}{}{‘Afrīn II, 35 n 63, 40 n. 139, 48 n. 220.}
\par\vspace{-0.09pt}
\hypertarget{AB01-PDF0808-col_left-E15}{}
\DLine{AB01-PDF0808-col_left-L025}{}{Afyūn Qarā Ḥiṣār II, 44 n. 183.}
\par\vspace{-0.17pt}
\hypertarget{AB01-PDF0808-col_left-E16}{}
\DLine{AB01-PDF0808-col_left-L026}{}{al-Aḥsā’ II, 52 n. 254.}
\par\vspace{-0.09pt}
\hypertarget{AB01-PDF0808-col_left-E17}{}
\DLine{AB01-PDF0808-col_left-L027}{}{al-Ahwāz II, \textit{35}.}
\par\vspace{0.18pt}
\hypertarget{AB01-PDF0808-col_left-E18}{}
\DLine{AB01-PDF0808-col_left-L028}{}{Akhlāṭ v. Khilāṭ.}
\end{minipage}
\hfill\vrule width.3pt\hfill\begin{minipage}[t]{.485\textwidth}\vspace{0pt}\fontsize{11.1}{14.02}\selectfont\setlength{\GutterOffset}{0pt}\setlength{\FullSourceWidth}{\linewidth}
\hypertarget{AB01-PDF0808-col_right-E01}{}
\DLine{AB01-PDF0808-col_right-L001}{}{Akhmīm II, \textit{52} (et 217).}
\par\vspace{0.01pt}
\hypertarget{AB01-PDF0808-col_right-E02}{}
\DLine{AB01-PDF0808-col_right-L002}{}{Akhsīkat II, 36 n. 85.}
\par\vspace{0.09pt}
\hypertarget{AB01-PDF0808-col_right-E03}{}
\DLine{AB01-PDF0808-col_right-L003}{}{‘Akkā II, \textit{39} (et 216).}
\par\vspace{-0.17pt}
\hypertarget{AB01-PDF0808-col_right-E04}{}
\DLine{AB01-PDF0808-col_right-L004}{}{Aksūm v. Auxume.}
\par\vspace{0.00pt}
\hypertarget{AB01-PDF0808-col_right-E05}{}
\DLine{AB01-PDF0808-col_right-L005}{}{Alani (\textit{al-Lān}) 172.}
\par\vspace{0.19pt}
\hypertarget{AB01-PDF0808-col_right-E06}{}
\DLine{AB01-PDF0808-col_right-L006}{}{Albania II, \textit{34}. — Albaniae Portae v. Porta.}
\par\vspace{0.18pt}
\hypertarget{AB01-PDF0808-col_right-E07}{}
\DLine{AB01-PDF0808-col_right-L007}{}{Albion Britannica (\textit{Alūyōn Brehṭanīqī}) II, \textit{33}}
\par\vspace{-0.36pt}
\DLine{AB01-PDF0808-col_right-L008}{}{\hspace*{10.04pt}nr. 2, 211.}
\par\vspace{0.01pt}
\hypertarget{AB01-PDF0808-col_right-E08}{}
\DLine{AB01-PDF0808-col_right-L009}{}{Alcaz v. Alqis.}
\par\vspace{0.18pt}
\hypertarget{AB01-PDF0808-col_right-E09}{}
\DLine{AB01-PDF0808-col_right-L010}{}{Aleppo v. Ḥalab.}
\par\vspace{0.19pt}
\hypertarget{AB01-PDF0808-col_right-E10}{}
\DLine{AB01-PDF0808-col_right-L011}{}{Alexandria (\textit{al-Iskandariyyah}) \textit{41} bis et n. 2,}
\par\vspace{0.10pt}
\DLine{AB01-PDF0808-col_right-L012}{}{\hspace*{10.04pt}\textit{42} et n. 3, 73 n. 3, 160, 172; II, \textit{38} (et}
\par\vspace{-0.09pt}
\DLine{AB01-PDF0808-col_right-L013}{}{\hspace*{10.04pt}216), 170 n. 2, 206, 210 plur., 212.}
\par\vspace{0.09pt}
\hypertarget{AB01-PDF0808-col_right-E11}{}
\DLine{AB01-PDF0808-col_right-L014}{}{Alexandria ad Issum v. Iskandarūnah.}
\par\vspace{-0.17pt}
\hypertarget{AB01-PDF0808-col_right-E12}{}
\DLine{AB01-PDF0808-col_right-L015}{}{Algeciras II, 220.}
\par\vspace{0.01pt}
\hypertarget{AB01-PDF0808-col_right-E13}{}
\DLine{AB01-PDF0808-col_right-L016}{}{Almeria v. al-Mariyyah.}
\par\vspace{0.18pt}
\hypertarget{AB01-PDF0808-col_right-E14}{}
\DLine{AB01-PDF0808-col_right-L017}{}{Alqis (?) II, \textit{50} et n.}
\par\vspace{-0.09pt}
\hypertarget{AB01-PDF0808-col_right-E15}{}
\DLine{AB01-PDF0808-col_right-L018}{}{‘Amawās II, \textit{47}.}
\par\vspace{0.01pt}
\hypertarget{AB01-PDF0808-col_right-E16}{}
\DLine{AB01-PDF0808-col_right-L019}{}{Āmid II, \textit{41}, 212.}
\par\vspace{0.36pt}
\hypertarget{AB01-PDF0808-col_right-E17}{}
\DLine{AB01-PDF0808-col_right-L020}{}{‘Ammūriyah (Amorium) II, \textit{44} (et 217), 212.}
\par\vspace{-0.26pt}
\hypertarget{AB01-PDF0808-col_right-E18}{}
\DLine{AB01-PDF0808-col_right-L021}{}{al-‘Amq II, \textit{35} et n.}
\par\vspace{0.01pt}
\hypertarget{AB01-PDF0808-col_right-E19}{}
\DLine{AB01-PDF0808-col_right-L022}{}{Āmul II, \textit{53}.}
\par\vspace{0.18pt}
\hypertarget{AB01-PDF0808-col_right-E20}{}
\DLine{AB01-PDF0808-col_right-L023}{}{‘Amwās v. ‘Amawās.}
\par\vspace{0.10pt}
\hypertarget{AB01-PDF0808-col_right-E21}{}
\DLine{AB01-PDF0808-col_right-L024}{}{Anafah II, 48 n. 221. Cfr. \textit{al-Mashriq}, V, 1902,}
\par\vspace{-0.27pt}
\DLine{AB01-PDF0808-col_right-L025}{}{\hspace*{10.04pt}p. 109-111.}
\par\vspace{0.46pt}
\hypertarget{AB01-PDF0808-col_right-E22}{}
\DLine{AB01-PDF0808-col_right-L026}{}{Anastasiopolis (\textit{Dārā}) II, 42 n. 159; (\textit{ar-Ru-}}
\par\vspace{-0.27pt}
\DLine{AB01-PDF0808-col_right-L027}{}{\hspace*{10.04pt}\textit{ṣāfah}) II, 45 n. 201.}
\par\vspace{0.27pt}
\hypertarget{AB01-PDF0808-col_right-E23}{}
\DLine{AB01-PDF0808-col_right-L028}{}{Anāwarzā II, 43 n. 179. Cfr. F. X. \Name{Schaffner},}
\end{minipage}
\end{document}
